\section{Introduction}

Let \(V\) be an \(n\)-dimensional vector space over a field \(k\). In this
paper we study the natural action of \(G=\GL(V)\) on
\[
  \Gr(2,\wedgepow{2}{V}).
\]
Equivalently, we study \(2\)-dimensional subspaces of \(\wedge^2 V\) up to the
action of \(\GL(V)\).

This problem lies inside the general finite-orbit theory for algebraic groups
on subspace varieties. Applied to the \(\GL(V)\)-module \(\wedge^2 V\), the
results of Guralnick, Liebeck, Macpherson, and Seitz \cite{GLMS1997} imply that
there are only finitely many \(\GL(V)\)-orbits on \(\Gr(2,\wedge^2 V)\) for
\(\dim V\le 7\).

There are two nearby pieces of literature. First, the weak-congruence
classification of alternating pencils due to Sergeichuk and to
Drozd--Plakosh provides the orbit types in principle
\cite{Sergeichuk2010,DrozdPlakosh2014,DrozdPlakosh2016}. Second, the papers of
Brooksbank--O'Brien, Brooksbank--Wilson, Brooksbank--Maglione--Wilson, and
Brooksbank--Li--Qiao--Wilson develop algorithms for preservers, isometry
groups, and pseudo-isometries of alternating matrix spaces
\cite{BrooksbankObrien2008,BrooksbankWilson2012,BrooksbankMaglioneWilson2017,BrooksbankLiQiaoWilson2020}.
What is still missing from this picture is a uniform exterior-square orbit
table through the finite range \(\dim V\le 7\): explicit wedge representatives,
their geometric stabilizers, and the resulting finite-field orbit tables. This
is the gap filled here.

For a representative \(L\le \wedge^2 V\), let
\[
  \rho_L:\Stab_{\GL(V)}(L)\to \GL(L)\cong \GL_2
\]
be the induced action on the \(2\)-space \(L\). The paper records
\[
  K_L=\ker(\rho_L),
  \qquad
  Q_L=\im(\rho_L),
\]
so that the stabilizer is described by the exact sequence
\[
  1\to K_L\to \Stab_{\GL(V)}(L)\to Q_L\to 1.
\]
The row-by-row geometric calculations establishing these kernel and quotient
descriptions are verified in a companion Lean~4 formalization
\cite{RizzoliWedge2Formalization2026}. The finite-field descent and orbit-count
identities are then checked in the paper, with companion Magma scripts used for
the symbolic orbit-count simplifications.

Weak-congruence labels are used only as names for the chosen representatives.
They are not used to determine the quotient \(Q_L\). Instead, \(Q_L\) is
computed directly as the image of \(\rho_L\). We keep row numbers for easy
navigation in the tables and in the Lean crosswalk.

In the stabilizer formulas, \(U_n\) denotes a connected \(n\)-dimensional
unipotent group, and \(U_n(q)\) denotes its group of \(\Fq\)-points. When the
unipotent radical is explicitly a vector group, we write \(\Ga^m\).

The main result is the following.

\begin{theorem}
\label{thm:intro-main}
Let \(k=\overline{\mathbf F}_q\), let \(V=k^n\), and assume \(n=4,5,6,7\).

\begin{enumerate}[label=\arabic*.]
\item The geometric \(\GL(V)\)-orbits on \(\Gr(2,\wedge^2 V)\) are represented
  by the weak-congruence rows listed in
  Section~\ref{sec:geometric-stabilizers}.
\item For each row \(L\), the algebraic stabilizer is determined by the exact
  sequence
  \[
    1\to K_L\to \Stab_{\GL(V)}(L)\to Q_L\to 1,
  \]
  with \(K_L\) and \(Q_L\) listed explicitly in
  Section~\ref{sec:geometric-stabilizers}. These geometric stabilizer
  calculations are verified row by row in the companion Lean~4 formalization.
\item Over \(\Fq\), the geometric families refine to exactly \(4,6,14,20\)
  finite-field orbits. Their stabilizers are listed in
  Section~\ref{sec:finite-fields}.
\item For each \(n=4,5,6,7\), the sum of the orbit sizes determined by those
  finite-field stabilizers equals the number of \(2\)-spaces in
  \(\wedge^2 \Fq^n\). Therefore the finite-field orbit tables are complete.
  Consequently the numbers of geometric orbits are \(3,5,11,16\).
\end{enumerate}
\end{theorem}

The companion repository contains the Lean~4 formalization of the geometric
stabilizer calculations and the Magma scripts used for the finite-field tables
and orbit-count identities. It is available at
\url{https://github.com/alunik/gl2spaces-wedge2}.

Section~\ref{sec:preliminaries} fixes notation and the weak-congruence naming
convention. Section~\ref{sec:geometric-stabilizers} gives the geometric orbit
and stabilizer tables together with the Lean crosswalk protocol.
Section~\ref{sec:finite-fields} carries out the finite-field descent, records
the finite-field orbit/stabilizer tables, proves the orbit-count identity, and
deduces the completeness of the geometric list.
